% Compact mode is enabled only by the English appendix layout.
\ifdefined\CompactOracleTable
\begin{table}[pos=!htbp]
\else
\begin{table*}[pos=!tbp]
\fi
\centering
\caption{Production-sensitivity alignment of the Ops labelling method used by PHGST-DM in the development pilot.}
\label{tab:oracle-pilot-ci}
\ifdefined\CompactOracleTable
\scriptsize
\setlength{\tabcolsep}{3pt}
\else
\footnotesize
\setlength{\tabcolsep}{8pt}
\fi
\renewcommand{\arraystretch}{1.1}
\begin{tabular}{@{}lrrr@{}}
\toprule
\ifdefined\CompactOracleTable
Metric & \makecell{Point\\estimate} & \makecell{95\% confidence\\interval} & \makecell{Interval\\width} \\
\else
Metric & Point estimate & 95\% confidence interval & Interval width \\
\fi
\midrule
NDCG@3 & 0.531 & [0.492, 0.570] & 0.079 \\
Positive Recall@2 & 55.556\% & [51.141, 59.971]\% & 8.830 pp \\
Positive Recall@3 & 68.519\% & [63.270, 73.768]\% & 10.497 pp \\
Weighted Recall@3 & 71.623\% & [66.052, 77.194]\% & 11.143 pp \\
Coverage@3$_{0.9}$ & 72.222\% & [66.974, 77.470]\% & 10.497 pp \\
\bottomrule
\end{tabular}
\par\vspace{3pt}
\begin{minipage}{\ifdefined\CompactOracleTable\columnwidth\else0.95\textwidth\fi}
\scriptsize
The pilot comprises 30 disturbance-free development episodes. Interval width is the upper bound minus the lower bound; pp denotes percentage points. These exploratory intervals do not establish superiority over the comparison methods or generalization to other disturbance conditions.
\end{minipage}
\ifdefined\CompactOracleTable
\end{table}
\else
\end{table*}
\fi
